\chapter{The method}

\section{Comparing texts by their byte pairs}

This chapter compares texts by how often each pair of adjacent bytes occurs in them. It counts the pairs, measures the distance between two texts, checks the measurement against two halves of the same text, corrects for how many pairs a text uses, groups the texts, and allows for languages that share an ancestor.

\subsection{1. Counting byte pairs}

Take a text \(T\), remove its whitespace, and encode it in UTF-8 as the bytes \(b_0 \dots b_{m-1}\). Each pair of neighboring bytes gets one index:

\(k_i \;=\; 256\,b_i + b_{i+1}, \qquad k_i \in [0,\,65536)\)

This turns a text of any length, written in 256 possible byte values, into a fixed array of \(2^{16}\) counts. Nothing in it decides what a character is. Every pair gets its own index. No two different pairs share a count.

\(P_T(k) \;=\; \frac{c_T(k)}{\sum_{j} c_T(j)}, \qquad c_T(k) = \bigl|\{\, i : k_i = k \,\}\bigr|\)

\(P_T\) is a probability distribution over \(2^{16}\) outcomes, the byte-pair histogram of the text.

\subsection{2. Total variation distance}

\(D(P,Q) \;=\; \tfrac{1}{2} \sum_{k=0}^{65535} \bigl| P(k) - Q(k) \bigr| \;\in\; [0,1]\)

This is the total variation distance. It is 0 when the two distributions agree everywhere and 1 when no outcome has nonzero probability in both.

\subsection{3. Split-half check}

A distance \(D(P_A,P_B)\) between two corpora means nothing until it is known how large a distance chance alone gives at that sample size. Split each corpus in half and measure the distance between its own halves:

\(D_{\text{self}}(T) \;=\; D\bigl(P_{T[0:m/2]},\; P_{T[m/2:m]}\bigr)\)

A distance between corpora can be read when

\(D_{\text{self}}(T) \;\ll\; D(P_T, P_{T'}) \quad \forall\, T' \neq T\)

and cannot be read when \(D_{\text{self}}(T) > \max_{T'} D_{\text{self}}(T')\), because then the corpus is too small to tell apart from noise.

On 20 corpora cut to 6707 bytes each, \(D_{\text{self}}\) ran from \(0.1118\) to \(0.3192\), with a mean of \(0.2361\).

\subsection{4. Correcting for how many pairs a text uses}

\(D_{\text{self}}\) falls as a distribution concentrates on fewer pairs, whatever the language. Report how many pairs occur, and the entropy, beside it:

\(\mathrm{supp}(P) = \bigl|\{\,k : P(k) > 0\,\}\bigr|, \qquad H(P) = -\sum_{k} P(k)\log_2 P(k)\)

Across the same 20 corpora, the number of pairs ran from \(193\) to \(970\) and rose and fell with \(D_{\text{self}}\) almost in step. A corpus in an encoding that uses few pairs looks consistent with itself for arithmetic reasons alone. Ranking writing systems by \(D_{\text{self}}\) ranks them by how many pairs they use, until shown otherwise.

\subsection{5. Text damaged by extraction}

Take the sets of characters \(S = \{s_1 \dots s_q\}\) that only a given writing system uses:

\(\mathrm{sets}(T) = \sum_{s \in S} \mathbf{1}\bigl[\,\exists\, c \in T,\ c \in s\,\bigr], \qquad \mu(T) = \frac{10^3}{|T_\alpha|}\bigl|\{\,c \in T : \mathrm{comb}(c)\,\}\bigr|\)

Here \(T_\alpha\) is the text's letters and \(\mathrm{comb}\) marks combining characters, such as accents written as separate code points. A damaged PDF extraction loses the combining marks first. A file with \(\mu \to 0\) and \(\mathrm{sets} \ll q\) has been damaged even though the text left in it still reads normally.

On 60 extracted papers with \(q=7\), one file scored \(\mathrm{sets} \ge 5\) with \(\mu > 0\). Every other file had lost the marks completely.

\subsection{6. Permutation test}

To tell a real split of the data from what splitting any set of keys produces by chance, shuffle the labels within each group the split started from, keeping the size of every cell the same \cite{src:Good-2005}:

\(\pi^{\ast} \sim \mathrm{Unif}\bigl(\mathcal{S}(\text{labels within group})\bigr), \qquad \Delta = \mathrm{E}\bigl[\,r(\pi^{\ast})\,\bigr] - r(\pi_{\text{obs}})\)

\(\Delta \le 0\) means the observed split gained nothing that a random split of the same shape would not. Over 19 corpora limited to 60000 tokens each, \(\Delta\) ran from \(+0.0056\) to \(+0.0901\), and the spread of the shuffled splits was at most \(0.0133\) over 5 draws.

\subsection{7. Checking the tree against the distances}

Group the texts by average-linkage clustering on \(D\). A tree can be built from any distance matrix, including one with no groups in it. Report the cophenetic correlation, the correlation between each measured distance and the height \(h(i,j)\) at which \(i\) and \(j\) first join in the tree \cite{src:Sokal-and-Rohlf-1962}:

\(\rho_c \;=\; \mathrm{corr}\bigl(\,D(i,j),\; h(i,j)\,\bigr)_{i < j}\)

A low \(\rho_c\) means the tree shows groups the distances do not support, and the result is only an ordering. On 19 languages, over \(171\) pairs, \(\rho_c = 0.8406\).

\subsection{8. Sampling one language per family}

Languages are not independent samples, because related languages share a history. Drawing one language from each family at a time removes that shared part:

\(\hat{\pi} \;=\; \frac{1}{d}\sum_{t=1}^{d} \frac{1}{|F|}\sum_{f \in F} \mathbf{1}\bigl[\,h(x_{f,t})\,\bigr], \qquad x_{f,t} \sim \mathrm{Unif}(f)\)

over \(d\) draws from the set of families \(F\). Report \(\hat{\pi}\) with its standard deviation across the draws. On 418 languages in 121 families with \(d=300\), the rates for each area ran from \(0.089\) to \(0.716\), with a spread of \(0.029\) to \(0.069\).
